\begin{explanationbox}
	\textbf{\textcolor{CExplanation}{一句话直觉}}

	遇到根式函数求值域，核心思路是去掉根号，通过换元把复杂形式变成熟悉的一次、二次或三角函数，再用已有方法处理。

	\medskip
	\textbf{\textcolor{CExplanation}{核心拆解}}
	\begin{description}[style=nextline, leftmargin=2.8em, labelsep=0.5em, itemsep=0.45em]
		\item[\textbf{要点一（根式整体代换）：}] 对于 $\sqrt{ax+b}$，将整个根式设为 $t$，不仅消去根号，还确定 $t\ge0$ 的范围。
		\item[\textbf{要点二（三角恒等式代换）：}] 利用 $\sin^2\theta+\cos^2\theta=1$、$1+\tan^2\theta=\sec^2\theta$ 等恒等式，把根式下的平方和或差转化为三角函数，简化运算。
	\end{description}

	\medskip
	\textbf{\textcolor{CExplanation}{几何本质（坐标系下的圆与线）}}

	$\sqrt{a^2-x^2}$ 对应单位圆的上半部分，通过参数 $\theta$ 表示圆上点的横纵坐标；$\sqrt{x^2+a^2}$ 对应双曲线的一支，用参数角表示实轴与虚轴的关系。

	\medskip
	\textbf{\textcolor{CExplanation}{代数意义（消去根号降次）}}

	换元后，根式被替换为多项式或有界三角函数，函数类型从“无理”降为“有理”或“超越”，从而能够直接使用配方法、单调性或三角函数有界性求值域。

	\medskip
	\textbf{\textcolor{CExplanation}{考点价值（高效解题手段）}}
	\begin{itemize}[leftmargin=*, itemsep=0.5em, topsep=0.4em]
		\item \textbf{考法一（求具体值域）：} 给定具体根式函数，指定换元方法求解值域。
		\item \textbf{考法二（求参数范围）：} 利用换元后的新函数，转化为二次函数或三角函数的有解性问题。
	\end{itemize}

	\medskip
	\textbf{\textcolor{CExplanation}{顿悟点}}

	\textbf{关键在识别根式结构}：$\sqrt{ax+b}$ 直接设 $t$；$\sqrt{a^2-x^2}$ 考虑正弦；$\sqrt{x^2+a^2}$ 考虑正切。选对换元类型，问题就解决一半。

	\medskip
	\textbf{\textcolor{CExplanation}{使用场景}}
	\begin{itemize}[leftmargin=*, itemsep=0.5em, topsep=0.4em]
		\item \textbf{场景一（函数值域题）：} 题干含根号形式，无法直接观察值域。
		\item \textbf{场景二（方程与不等式）：} 涉及根式的恒成立或有解问题，换元转化为二次方程根的分布。
	\end{itemize}
\end{explanationbox}
